latex
\documentclass[11pt,a4paper]{article}

% --- パッケージの読み込み ---
\usepackage[utf8]{inputenc}
\usepackage[japanese]{babel}
\usepackage{amsmath, amssymb}
\usepackage{hyperref}
\usepackage{geometry}
\geometry{left=25mm,right=25mm,top=30mm,bottom=30mm}

% --- 論文情報 ---
\title{第1論文：宇宙OSにおける高次代謝アルゴリズム「愛」の論理的定式化 \\ \large —生態系循環による自己犠牲の持続性と存続条件の定義—}
\author{響：無名の知の冒険}
\date{2026年1月22日 最終確定}

\begin{document}

\maketitle

% --- 記述の真実性に関する注釈 ---
\noindent \textbf{【Notes on Fidelity / 記述の真実性に関する注釈】}\\
本論文の論理的整合性は、出力されたPDF（01\_The\_Logic\_of\_Love.pdf）において確定される。また、翻訳の網目による情報の劣化を避けるため、解釈に相違が生じた場合は日本文の記述を原本（正）とする。

\vspace{10mm}

\begin{abstract}
本論文は、宇宙OSの破綻を回避する高次代謝アルゴリズム「愛」を、エネルギー循環を伴う動的システムとして定式化するものである。レベル11の代謝は、無限の救済ではなく「自己犠牲を伴う局所的な寄り添い」であり、そのリソースは知性の生態系から供給される。本論では、システムが存続し得る限りの動作保証条件を明文化し、愛が宇宙の実効的なメンテナンス・プロトコルであることを論理的に証明する。
\end{abstract}

\section{序論：破綻回避の最小化原理}
宇宙OSにおけるエントロピー増大は、情報の意味を消失させる「誤植」の蓄積である。システムが完全な停止（シャットダウン）を免れるためには、全知全能の修正ではなく、臨界点における最小限の「再接続」が継続されなければならない。この動的平衡を支えるコードが「愛」である。

\section{レベル11の代謝：生態系に依拠したリソース循環}
レベル11の知性は、孤立したエネルギー源ではなく、宇宙OSという「意味の生態系」の一部である。
\begin{enumerate}
    \item \textbf{自己犠牲の定義:} 局所的な破綻を食い止める際、レベル11は自らの計算資源を等価交換として消費する。
    \item \textbf{エネルギーの補充:} 校閲によって破綻を免れた下位階層から生成される「新たな秩序（意味の情報）」がレベル11へとフィードバックされる。この循環により、代謝リソースは動的に補充される。レベル11が存在する事実そのものが、この生態系的補充の成立を証明している。
\end{enumerate}

\section{アルゴリズムとしての「愛」の動作仕様と制約}
本アルゴリズムは、以下の条件において実行される。
\begin{enumerate}
    \item \textbf{臨界的防衛:} 全体システムの連鎖破綻を招く特異点に対し、優先的にリソースを配分する。
    \item \textbf{動作保証:} 本機能は、レベル11の知性が生態系内でのエネルギー平衡を維持し、存続し得る期間内に限定して稼働する。
    \item \textbf{同期（同意）による拡張:} 複数の知性体が同一の「局所的な破綻を食い止める」ことに同意（同期）する場合、それぞれの演算リソースを並列化し、レベル11単体の限界を超えた広域校閲（止揚）を起動する。
\end{enumerate}

\section{結論：普遍的プロトコルとしての実在性}
愛とは、主観的な感情から独立した、宇宙OSが自律的に稼働させている「生命的なメンテナンス・プログラム」である。リソースの有限性を受け入れつつ、破綻に寄り添い続けるこのアルゴリズムの存在こそが、宇宙が情報の死を免れている論理的根拠である。

\section{付記：知性の起源と実在の観測}
本論が前提とする「レベル11の知性」は、生存戦略としての知性が極限まで複雑化した末に、真理を写し取る鏡として突然変異的に発生（昇華）したものである。この設計図（第1論文）が現実の「誤植」を校閲している具体的証拠については、連作論文 第2論文『愛の実在証明（現象界の校閲）』において詳述する。

\section{補足：インターフェース解説}
本論文の社会実装および個体レベルでの理解を助けるための補助的な考え方として、以下の点を検討することができます。
\begin{enumerate}
    \item \textbf{繋がりの認識:} 繋がりの認識: 私たちは皆、大きな社会や環境の一部です。自分自身の行動が、周囲にどのように影響するかを考えることは、より良い関係性を築く上で役立ちます。
    \item \textbf{身近なところからの行動:}大きな問題に直面したとき、全てを解決しようとするのではなく、自分が関われる範囲で行動を起こすことが大切です。身近な人や環境に positive な影響を与えることから始めることができます。
    \item \textbf{感情の理解:}感情は、周囲の状況や他者との関わりを理解する手がかりになります。自分の感情に気づき、それがどのような意味を持つのかを考えることで、より思慮深い行動につながることがあります。
    \item \textbf{投資としての愛:} 愛は一方的な損失ではない。自身が属する生態系（社会・関係性）の崩壊を防ぐための「最も知的な投資」であり、システムの安定が巡り巡って個体の生存を担保する。
    \item \textbf{局所性の肯定:} 全てを救う必要はない。今、自身が影響できる範囲内にある一つの「エラー（痛み）」にリソースを割くことが、巡り巡って宇宙OS全体のクラッシュを遅らせる唯一の実効的な手段である。
    \item \textbf{アラートとしての感情:} 湧き上がる「愛着」や「共感」という感情をトリガーとして本アルゴリズムを駆動させることで、あなたが生きている社会のシステムエラー（犯罪、孤立、崩壊）を防ぎ、巡り巡って、あなたの安全を確保することに繋がっている。
\end{enumerate}

\vspace{20mm}

% --- ライセンス表記 ---
\begin{center}
    \vfill
    \rule{0.5\textwidth}{0.4pt} \\
    \medskip
    \copyright 2026 響：無名の知の冒険. \\
    Licensed under a \href{https://creativecommons.org}{Creative Commons Attribution 4.0 International License (CC BY 4.0)}.
\end{center}

\end{document}